
\subsection{Théorie} \label{sec:theorie}
À partir de la Ref. \cite{protocoleDMA}, nous pouvons décrire, depuis les lois de Hooke et de Newton, la modélisation de notre système. L’instrument impose une contrainte périodique en fonction du temps $\sigma$(t) de pulsation $\omega$ et d’une contrainte de référence $\sigma_0$ telle que
\begin{equation}
    \sigma(t) = \sigma_0 \sin(\omega t).
	\label{eq:1}
\end{equation}

Notons une déformation $\epsilon$ sous déphasage d’un angle $\delta$ :
\begin{equation}
    \varepsilon(t) = \varepsilon_0 \sin(\omega t + \delta).
	\label{eq:2}
\end{equation}

Selon les hypothèses du protocole (Ref. \cite{protocoleDMA} p. 2-3), les grandeurs viscoélastiques mesurées $E′$ (module de stockage) et $E′′$ (module de perte) sont, après dérivations des Eqs. \ref{eq:1} et \ref{eq:2} :
\begin{equation}
	E' &= \frac{\sigma_0}{\varepsilon_0}\cos\delta \hspace{0.4cm} et \hspace{0.4cm}
    E'' &= \frac{\sigma_0}{\varepsilon_0}\sin\delta.
	\label{eq:3}
\end{equation}

Elles permettent d’écrire le pic du facteur d’amortissement $\tan \delta$ correspondant à la transition vitreuse mécanique à partir de relations trigonométriques avec les Eqs. \ref{eq:3}, tel que :
\begin{equation}
	\tan\delta &= \frac{E'}{E'\'}.
	\label{eq:4}
\end{equation}
Comme cette transition dépend de la fréquence d’excitation, les températures des pics mesurés sont exploitées à l’aide de la loi d’Arrhenius (voir Ref. \cite{protocoleDMA}) :
\begin{equation}
    f = f_0 \exp\!\left(-\frac{E_a}{RT}\right),
	\label{eq:arrhenius}
\end{equation}

que l’on linéarise sous la forme

\begin{equation}
    \ln f = \ln f_0 - \frac{E_a}{R}\,\frac{1}{T}.
	\label{eq:arrheniuslin}
\end{equation}

Ainsi, la représentation de $\ln f$ en fonction de $1/T_g(f)$ permet d’extraire l’énergie
d’activation $E_a$ associée à la transition vitreuse du PMMA.

